\section{Résultats}

\subsection{Oscillations libres}

Avant de mesurer une oscillation, nous avons mesure le rapport de force qu'il faut exercer sur le balancier afin de le placer a un certain angle.
\\
\noindent \begin{minipage}{0.5\textwidth}
\includegraphics[width=\textwidth]{raideur.png}
	\captionof{figure}{$M_r$($\theta$)}
\label{fig:C}
\end{minipage}
\begin{minipage}{0.4\textwidth}
	Nous obtenons par consequent le graphique (\ref{fig:C}) ci-joint.
	Grace a sa pente d'unite ... nous avons pu directement extraire la valeur C, soit la constante de raideur du ressort utilise grace a la relation (\ref{eq:nature}).\\ Soit :
	\\      C = 0.0123 +/- 0.0002 [Nm / $rad^-1$]
\noindent \end{minipage}
\\
\\
\\
En un second temps, apres avoir montee l'installation pour mesures les oscillations du balencier-spiral (voir fig:\ref{fig:samplesetup}), nous avons donne une position hors-equilibre au balencier et mesurer en fonction du frein imposee a les differences oscillations observable en annexe (voir fig:\ref{fig:oscillations}).

\begin{minipage}[t]{0.45\textwidth}
    \centering
    \includegraphics[width=\linewidth]{lambda_quadratique.png}
	\captionof{figure}{$\lambda$(I)}
	\label{fig:lambda-current_quad}
\end{minipage}
\hfill
\begin{minipage}[t]{0.5\textwidth}
    \centering
    \includegraphics[width=\linewidth]{lambda.png}
	\captionof{figure}{$\lambda$($I^2$)}
	\label{fig:lambda-current}
\end{minipage}
\\
\\
La relation entre le facteur d'amortissement et le courant circulant dans les freins n'est pas lineaire. Nous constatons verifions la relation (\ref{eq:lambda-courant}), avec le graphique \ref{fig:lambda-current}. Nous pouvons extraire grace a la relation \ref{eq:JCW} la moyenne des pulsations propres $\overline{\omega_0}$ = 3.25 +/- 0.02 [rad/s] (voir tab:\ref{val:omega0}) desquelles on extrait la frequence propre : $f_0$ = ... (a partir de l'inverse de la periode : Eq. \ref{eq:t}). Avec la constante de torsion du ressort le moment d'inertie du balencier $J$ = $C$ / $\omega_^2$ = $0.0123/3.25^2$ = (1.16 +/- 0.000) \cdot $10^3$  [kg / $m^2$]. 

\newpage
\subsection{Expérience 2}

En tracant la phase en fonction de la frequence (voir \ref{fig:phase}), nous pouvons retrouver les frequences propres $f_0$ pour les alimentations a 0.5 en bleu et 0.9A en rouge.
        Les frequences propres resultent du point d'inflexions de leurs courbes correspondantes.
\\
\noindent \begin{minipage}{0.65\textwidth}
\vspace{0.4cm}
\includegraphics[width=\textwidth]{force_phase.png}
	\captionof{figure}{$\phi$(f), evolution de la Phase}
\label{fig:phase}
\end{minipage}
\hspace{0.00005\textwidth}
\begin{minipage}{0.4\textwidth}
	Ci-joint (\ref{fig:phase}), les frequences propres ont pu s'obtenir a partir de la relation de la phase-freqence (voir \ref{eq:phase_forcee}), sachant que la phase en fonction de la pulsation propore vaut $\pi/2$. L'ajustement nous donne :
	\\ \\ $f_0$(0.5) = 0.518 $+/-$ 0.009 [Hz].
	\\ \\ $f_0$(0.9) = 0.528 $+/-$ 0.004 [Hz].
\noindent \end{minipage}

\noindent \begin{minipage}{0.65\textwidth}
\vspace{0.4cm}
\includegraphics[width=\textwidth]{force_amp.png}
	\captionof{figure}{A(f), evolution de l'elongation angulaire}
	\label{fig:amplitude}
\label{fig:nature}
\end{minipage}
\hspace{0.00005\textwidth}
\begin{minipage}{0.4\textwidth}
	En revanche, l'amplitude du balencier en fonction de la frequence nous permet d'obtenir les frequences propres en fonction de Eq. \ref{eq:omega_res}, soit en explorant le lien entre pulsation propre et pseudo pulsation :
	\\ \\ $f_0$(0.5) = 0.518 +/- 0.004 [Hz].
	\\ \\ $f_0$(0.9) = 0.519 +/- 0.001 [Hz]
	
\noindent \end{minipage}
\\

En fonction du graph d'amplitude,
nous obtenons en realisant un fit a partir du modele Eq. \ref{eq:amplitude_forcee} les valeurs suivantes :
\\\\ 
\textbf{$\lambda$(0.5) = 0.191 ± 0.003 [$1/s$] , \hspace{2.5cm} $\lambda$(0.9) = 0.609 ± 0.005 [$1/s$]. \\$\omega_0$(0.5) = 3.249 ± 0.001 [rad/s],  \hspace{2cm} $\omega_0$(0.9) = 3.315 ± 0.003 [rad/s]} \\\\Et la par moyenne des freqences obtenues : 
\\- \hspace{5cm} \textbf{$\overline{f_0}$ = 0.520 +/- 0.005 [Hz].} 

\vspace{0.5cm}
A partir de la puissance relative, si nous etablissons donc le rapport $\overline{P}$ / $P_0$ en fonction de la frequence sous differents courants, nous pouvons observer une courbe plus aplatie (voir \ref{fig:power9}) qu'en le graphique a 0.5 A (voir \ref{fig:power5}).
Ces courbes vont nous permettre d'extraire FMWH grace a des outils MATLAB et les facteurs de qualite Q a partir de la relation  Eq. \ref{eq:Q_fwhm}.
\\
\begin{minipage}[t]{0.5\textwidth}
    \centering
	\includegraphics[scale=0.3]{power_0.5A.png}
	\captionof{figure}{$P$/$P_0$(f) 0.5A}
	\label{fig:power5}
\end{minipage}
\hfill
\begin{minipage}[t]{0.5\textwidth}
    \centering
	\includegraphics[scale=0.3]{power_0.9A.png}
	\captionof{figure}{$P$/$P_0$(f) a 0.9A}
	\label{fig:power9}
\end{minipage}

\vspace{1cm}
À partir du rapport de puissance, nous obtenons la FWHM, directement liée à la largeur de la résonance :

\[
\mathrm{FWHM}_{0.5\,A} = 0.133998 \pm 0.003146~\mathrm{Hz}, \qquad
\mathrm{FWHM}_{0.9\,A} = 0.214367 \pm 0.002802~\mathrm{Hz}.
\]

Ce qui nous donne respectivement :
\[
	\mathrm{Q}_{0.5\A} = 1\qquad
	\mathrm{Q}_{0.9\A} = 3
\]
